\iftestbuild % TEST-ONLY-BEGIN
\setcounter{section}{3}\setcounter{theorem}{67}
\fi % TEST-ONLY-END
% Source: book pp. 370--382 (PDF 384--396); no solutions.
% \exnote: Axler's small italic comment under an exercise
% Numbered equations (9.75, 9.81): \axeqstep{ax:N.M} steps Axler's counter and
% sets the label just before the display; \axeqtag prints the number.
\DeclareRobustCommand\axeqstep[1]{\refstepcounter{theorem}\label{#1}}
\DeclareRobustCommand\axeqtag{\tag*{\thetheorem}}
\section{Tensor Products}

\subsection{Tensor Product of Two Vector Spaces}

The motivation for our next topic comes from wanting to form the product of a
vector $v\in V$ and a vector $w\in W$. This product will be denoted by
$v\otimes w$, pronounced ``$v$ tensor $w$'', and will be an element of some new
vector space called $V\otimes W$ (also pronounced ``$V$ tensor $W$'').

We already have a vector space $V\times W$ (see Section~3E), called the product
of $V$ and $W$. However, $V\times W$ will not serve our purposes here because it
does not provide a natural way to multiply an element of $V$ by an element of
$W$. We would like our tensor product to satisfy some of the usual properties
of multiplication. For example, we would like the distributive property to be
satisfied, meaning that if $v_1,v_2,v\in V$ and $w_1,w_2,w\in W$, then
\[ (v_1+v_2)\otimes w = v_1\otimes w+v_2\otimes w \quad\text{and}\quad
   v\otimes(w_1+w_2) = v\otimes w_1+v\otimes w_2. \]

\marginnote{To produce $\otimes$ in \TeX, type \texttt{\textbackslash otimes}.}%
We would also like scalar multiplication to interact well with this new
multiplication, meaning that
\[ \lambda(v\otimes w) = (\lambda v)\otimes w = v\otimes(\lambda w) \]
for all $\lambda\in\F$, $v\in V$, and $w\in W$.

Furthermore, it would be nice if each basis of $V$ when combined with each
basis of $W$ produced a basis of $V\otimes W$. Specifically, if
$e_1,\dots,e_m$ is a basis of $V$ and $f_1,\dots,f_n$ is a basis of $W$, then
we would like a list (in any order) consisting of $e_j\otimes f_k$, as $j$
ranges from $1$ to $m$ and $k$ ranges from $1$ to $n$, to be a basis of
$V\otimes W$. This implies that $\dim(V\otimes W)$ should equal
$(\dim V)(\dim W)$. Recall that $\dim(V\times W)=\dim V+\dim W$ (see
\ref{ax:3.92}), which shows that the product $V\times W$ will not serve our
purposes here.

To produce a vector space whose dimension is $(\dim V)(\dim W)$ in a natural
fashion from $V$ and $W$, we look at the vector space of bilinear functionals,
as defined below.

\begin{definition}[Bilinear functional on $V\times W$, the vector space
$\mathcal{B}(V,W)$]\label{ax:9.68}\leavevmode
\begin{itemize}
\item A \emph{bilinear functional} on $V\times W$ is a function
  $\beta\colon V\times W\to\F$ such that $v\mapsto\beta(v,w)$ is a linear
  functional on $V$ for each $w\in W$ and $w\mapsto\beta(v,w)$ is a linear
  functional on $W$ for each $v\in V$.
\item The vector space of bilinear functionals on $V\times W$ is denoted by
  $\mathcal{B}(V,W)$.
\end{itemize}
\end{definition}

If $W=V$, then a bilinear functional on $V\times W$ is a bilinear form; see
\ref{ax:9.1}.

The operations of addition and scalar multiplication on $\mathcal{B}(V,W)$ are
defined to be the usual operations of addition and scalar multiplication of
functions. As you can verify, these operations make $\mathcal{B}(V,W)$ into a
vector space whose additive identity is the zero function from $V\times W$ to
$\F$.

\begin{example}[Bilinear functionals]\label{ax:9.69}\leavevmode
\begin{itemize}
\item Suppose $\varphi\in V'$ and $\tau\in W'$. Define
  $\beta\colon V\times W\to\F$ by $\beta(v,w)=\varphi(v)\tau(w)$. Then $\beta$
  is a bilinear functional on $V\times W$.
\item Suppose $v\in V$ and $w\in W$. Define $\beta\colon V'\times W'\to\F$ by
  $\beta(\varphi,\tau)=\varphi(v)\tau(w)$. Then $\beta$ is a bilinear
  functional on $V'\times W'$.
\item Define $\beta\colon V\times V'\to\F$ by $\beta(v,\varphi)=\varphi(v)$.
  Then $\beta$ is a bilinear functional on $V\times V'$.
\item Suppose $\varphi\in V'$. Define $\beta\colon V\times\Lin(V)\to\F$ by
  $\beta(v,T)=\varphi(Tv)$. Then $\beta$ is a bilinear functional on
  $V\times\Lin(V)$.
\item Suppose $m$ and $n$ are positive integers. Define
  $\beta\colon\F^{m,n}\times\F^{n,m}\to\F$ by $\beta(A,B)=\tr(AB)$. Then
  $\beta$ is a bilinear functional on $\F^{m,n}\times\F^{n,m}$.
\end{itemize}
\end{example}

\begin{result}[Dimension of the vector space of bilinear
functionals]\label{ax:9.70}\leavevmode
\[ \dim\mathcal{B}(V,W)=(\dim V)(\dim W). \]
\end{result}

\begin{proof}
Let $e_1,\dots,e_m$ be a basis of $V$ and $f_1,\dots,f_n$ be a basis of $W$.
For a bilinear functional $\beta\in\mathcal{B}(V,W)$, let $\Mat(\beta)$ be the
$m$-by-$n$ matrix whose entry in row $j$, column $k$ is $\beta(e_j,f_k)$. The
map $\beta\mapsto\Mat(\beta)$ is a linear map of $\mathcal{B}(V,W)$ into
$\F^{m,n}$.

For a matrix $C\in\F^{m,n}$, define a bilinear functional $\beta_C$ on
$V\times W$ by
\[ \beta_C(a_1e_1+\cdots+a_me_m,\,b_1f_1+\cdots+b_nf_n)
   = \sum_{k=1}^n\sum_{j=1}^m C_{j,k}a_jb_k \]
for $a_1,\dots,a_m,b_1,\dots,b_n\in\F$.

The linear map $\beta\mapsto\Mat(\beta)$ from $\mathcal{B}(V,W)$ to $\F^{m,n}$
and the linear map $C\mapsto\beta_C$ from $\F^{m,n}$ to $\mathcal{B}(V,W)$ are
inverses of each other because $\beta_{\Mat(\beta)}=\beta$ for all
$\beta\in\mathcal{B}(V,W)$ and $\Mat(\beta_C)=C$ for all $C\in\F^{m,n}$, as you
should verify.

Thus both maps are isomorphisms and the two spaces that they connect have the
same dimension. Hence
$\dim\mathcal{B}(V,W)=\dim\F^{m,n}=mn=(\dim V)(\dim W)$.
\end{proof}

Several different definitions of $V\otimes W$ appear in the mathematical
literature. These definitions are equivalent to each other, at least in the
finite-dimensional context, because any two vector spaces of the same dimension
are isomorphic.

The result above states that $\mathcal{B}(V,W)$ has the dimension that we seek,
as do $\Lin(V,W)$ and $\F^{\dim V,\dim W}$. Thus it may be tempting to define
$V\otimes W$ to be $\mathcal{B}(V,W)$ or $\Lin(V,W)$ or $\F^{\dim V,\dim W}$.
However, none of those definitions would lead to a basis-free definition of
$v\otimes w$ for $v\in V$ and $w\in W$.

The following definition, while it may seem a bit strange and abstract at
first, has the huge advantage that it defines $v\otimes w$ in a basis-free
fashion. We define $V\otimes W$ to be the vector space of bilinear functionals
on $V'\times W'$ instead of the more tempting choice of the vector space of
bilinear functionals on $V\times W$.

\begin{definition}[Tensor product, $V\otimes W$, $v\otimes w$]\label{ax:9.71}\leavevmode
\begin{itemize}
\item The \emph{tensor product} $V\otimes W$ is defined to be
  $\mathcal{B}(V',W')$.
\item For $v\in V$ and $w\in W$, the \emph{tensor product} $v\otimes w$ is the
  element of $V\otimes W$ defined by
  \[ (v\otimes w)(\varphi,\tau) = \varphi(v)\tau(w) \]
  for all $(\varphi,\tau)\in V'\times W'$.
\end{itemize}
\end{definition}

We can quickly prove that the definition of $V\otimes W$ gives it the desired
dimension.

\begin{result}[Dimension of the tensor product of two vector
spaces]\label{ax:9.72}\leavevmode
\[ \dim(V\otimes W)=(\dim V)(\dim W). \]
\end{result}

\begin{proof}
Because a vector space and its dual have the same dimension (by
\ref{ax:3.111}), we have $\dim V'=\dim V$ and $\dim W'=\dim W$. Thus
\ref{ax:9.70} tells us that the dimension of $\mathcal{B}(V',W')$ equals
$(\dim V)(\dim W)$.
\end{proof}

To understand the definition of the tensor product $v\otimes w$ of two vectors
$v\in V$ and $w\in W$, focus on the kind of object it is. An element of
$V\otimes W$ is a bilinear functional on $V'\times W'$, and in particular it is
a function from $V'\times W'$ to $\F$. Thus for each element of $V'\times W'$,
it should produce an element of $\F$. The definition above has this behavior,
because $v\otimes w$ applied to a typical element $(\varphi,\tau)$ of
$V'\times W'$ produces the number $\varphi(v)\tau(w)$.

The somewhat abstract nature of $v\otimes w$ should not matter. The important
point is the behavior of these objects. The next result shows that tensor
products of vectors have the desired bilinearity properties.

\begin{result}[Bilinearity of tensor product]\label{ax:9.73}
Suppose $v,v_1,v_2\in V$ and $w,w_1,w_2\in W$ and $\lambda\in\F$. Then
\[ (v_1+v_2)\otimes w = v_1\otimes w+v_2\otimes w \quad\text{and}\quad
   v\otimes(w_1+w_2) = v\otimes w_1+v\otimes w_2 \]
and
\[ \lambda(v\otimes w) = (\lambda v)\otimes w = v\otimes(\lambda w). \]
\end{result}

\begin{proof}
Suppose $(\varphi,\tau)\in V'\times W'$. Then
\begin{align*}
\bigl((v_1+v_2)\otimes w\bigr)(\varphi,\tau) &= \varphi(v_1+v_2)\tau(w)\\
  &= \varphi(v_1)\tau(w)+\varphi(v_2)\tau(w)\\
  &= (v_1\otimes w)(\varphi,\tau)+(v_2\otimes w)(\varphi,\tau)\\
  &= (v_1\otimes w+v_2\otimes w)(\varphi,\tau).
\end{align*}
Thus $(v_1+v_2)\otimes w=v_1\otimes w+v_2\otimes w$.

The other two equalities are proved similarly.
\end{proof}

Lists are, by definition, ordered. The order matters when, for example, we
form the matrix of an operator with respect to a basis. For lists in this
section with two indices, such as
$\{e_j\otimes f_k\}_{j=1,\dots,m;\,k=1,\dots,n}$ in the next result, the
ordering does not matter and we do not specify it---just choose any convenient
ordering.

The linear independence of elements of $V\otimes W$ in (a) of the result below
captures the idea that there are no relationships among vectors in
$V\otimes W$ other than the relationships that come from bilinearity of the
tensor product (see \ref{ax:9.73}) and the relationships that may be present
due to linear dependence of a list of vectors in $V$ or a list of vectors in
$W$.

\begin{result}[Basis of $V\otimes W$]\label{ax:9.74}
Suppose $e_1,\dots,e_m$ is a list of vectors in $V$ and $f_1,\dots,f_n$ is a
list of vectors in $W$.
\begin{parts}
\item If $e_1,\dots,e_m$ and $f_1,\dots,f_n$ are both linearly independent
  lists, then
  \[ \{e_j\otimes f_k\}_{j=1,\dots,m;\,k=1,\dots,n} \]
  is a linearly independent list in $V\otimes W$.
\item If $e_1,\dots,e_m$ is a basis of $V$ and $f_1,\dots,f_n$ is a basis of
  $W$, then the list $\{e_j\otimes f_k\}_{j=1,\dots,m;\,k=1,\dots,n}$ is a
  basis of $V\otimes W$.
\end{parts}
\end{result}

\begin{proof}
To prove (a), suppose $e_1,\dots,e_m$ and $f_1,\dots,f_n$ are both linearly
independent lists. This linear independence and the linear map lemma
(\ref{ax:3.4}) imply that there exist $\varphi_1,\dots,\allowbreak\varphi_m\in V'$ and
$\tau_1,\dots,\tau_n\in W'$ such that
\[ \varphi_j(e_k) = \begin{cases} 1 & \text{if } j=k,\\
                                  0 & \text{if } j\ne k \end{cases}
   \quad\text{and}\quad
   \tau_j(f_k) = \begin{cases} 1 & \text{if } j=k,\\
                               0 & \text{if } j\ne k, \end{cases} \]
where $j,k\in\{1,\dots,m\}$ in the first equation and $j,k\in\{1,\dots,n\}$ in
the second equation.

Suppose $\{a_{j,k}\}_{j=1,\dots,m;\,k=1,\dots,n}$ is a list of scalars such
that\axeqstep{ax:9.75}
\begin{equation*}
\sum_{k=1}^n\sum_{j=1}^m a_{j,k}(e_j\otimes f_k) = 0. \axeqtag
\end{equation*}
Note that $(e_j\otimes f_k)(\varphi_M,\tau_N)$ equals $1$ if $j=M$ and $k=N$,
and equals $0$ otherwise. Thus applying both sides of \ref{ax:9.75} to
$(\varphi_M,\tau_N)$ shows that $a_{M,N}=0$, proving that
$\{e_j\otimes f_k\}_{j=1,\dots,m;\,k=1,\dots,n}$ is linearly independent.

Now (b) follows from (a), the equation $\dim V\otimes W=(\dim V)(\dim W)$ [see
\ref{ax:9.72}], and the result that a linearly independent list of the right
length is a basis (see \ref{ax:2.38}).
\end{proof}

Every element of $V\otimes W$ is a finite sum of elements of the form
$v\otimes w$, where $v\in V$ and $w\in W$, as implied by (b) in the result
above. However, if $\dim V>1$ and $\dim W>1$, then Exercise~4 shows that
\[ \{v\otimes w : (v,w)\in V\times W\} \ne V\otimes W. \]

\begin{example}[Tensor product of element of $\F^m$ with element of
$\F^n$]\label{ax:9.76}
Suppose $m$ and $n$ are positive integers. Let $e_1,\dots,e_m$ denote the
standard basis of $\F^m$ and let $f_1,\dots,f_n$ denote the standard basis of
$\F^n$. Suppose
\[ v = (v_1,\dots,v_m)\in\F^m \quad\text{and}\quad w = (w_1,\dots,w_n)\in\F^n. \]
Then
\begin{align*}
v\otimes w &= \biggl(\sum_{j=1}^m v_je_j\biggr)\otimes
              \biggl(\sum_{k=1}^n w_kf_k\biggr)\\
           &= \sum_{k=1}^n\sum_{j=1}^m (v_jw_k)(e_j\otimes f_k).
\end{align*}
Thus with respect to the basis
$\{e_j\otimes f_k\}_{j=1,\dots,m;\,k=1,\dots,n}$ of $\F^m\otimes\F^n$ provided
by \ref{ax:9.74}(b), the coefficients of $v\otimes w$ are the numbers
$\{v_jw_k\}_{j=1,\dots,m;\,k=1,\dots,n}$. If instead of writing these numbers
in a list, we write them in an $m$-by-$n$ matrix with $v_jw_k$ in row $j$,
column $k$, then we can identify $v\otimes w$ with the $m$-by-$n$ matrix
\[ \begin{pmatrix} v_1w_1 & \cdots & v_1w_n\\
                          & \ddots &       \\
                   v_mw_1 & \cdots & v_mw_n \end{pmatrix}. \]
\end{example}

See Exercises 5 and 6 for practice in using the identification from the
example above.

We now define bilinear maps, which differ from bilinear functionals in that the
target space can be an arbitrary vector space rather than just the scalar
field.

\begin{definition}[Bilinear map]\label{ax:9.77}\emergencystretch=1.5em
A \emph{bilinear map} from $V\times W$ to a vector space $U$ is a function
$\Gamma\colon V\times W\to U$ such that $v\mapsto\Gamma(v,w)$ is a linear map
from $V$ to $U$ for each $w\in W$ and $w\mapsto\Gamma(v,w)$ is a linear map
from $W$ to $U$ for each $v\in V$.
\end{definition}

\begin{example}[Bilinear maps]\label{ax:9.78}\leavevmode
\begin{itemize}
\item Every bilinear functional on $V\times W$ is a bilinear map from
  $V\times W$ to $\F$.
\item The function $\Gamma\colon V\times W\to V\otimes W$ defined by
  $\Gamma(v,w)=v\otimes w$ is a bilinear map from $V\times W$ to $V\otimes W$
  (by \ref{ax:9.73}).
\item The function $\Gamma\colon\Lin(V)\times\Lin(V)\to\Lin(V)$ defined by
  $\Gamma(S,T)=ST$ is a bilinear map from $\Lin(V)\times\Lin(V)$ to $\Lin(V)$.
\item The function $\Gamma\colon V\times\Lin(V,W)\to W$ defined by
  $\Gamma(v,T)=Tv$ is a bilinear map from $V\times\Lin(V,W)$ to $W$.
\end{itemize}
\end{example}

Tensor products allow us to convert bilinear maps on $V\times W$ into linear
maps on $V\otimes W$ (and vice versa), as shown by the next result. In the
mathematical literature, (a) of the result below is called the ``universal
property'' of tensor products.

\begin{result}[Converting bilinear maps to linear maps]\label{ax:9.79}
Suppose $U$ is a vector space.
\begin{parts}
\item Suppose $\Gamma\colon V\times W\to U$ is a bilinear map. Then there
  exists a unique linear map $\hat\Gamma\colon V\otimes W\to U$ such that
  \[ \hat\Gamma(v\otimes w) = \Gamma(v,w) \]
  for all $(v,w)\in V\times W$.
\item Conversely, suppose $T\colon V\otimes W\to U$ is a linear map. Then there
  exists a unique bilinear map $T^\#\colon V\times W\to U$ such that
  \[ T^\#(v,w) = T(v\otimes w) \]
  for all $(v,w)\in V\times W$.
\end{parts}
\end{result}

\begin{proof}
Let $e_1,\dots,e_m$ be a basis of $V$ and let $f_1,\dots,f_n$ be a basis of
$W$. By the linear map lemma (\ref{ax:3.4}) and \ref{ax:9.74}(b), there exists
a unique linear map $\hat\Gamma\colon V\otimes W\to U$ such that
\[ \hat\Gamma(e_j\otimes f_k) = \Gamma(e_j,f_k) \]
for all $j\in\{1,\dots,m\}$ and $k\in\{1,\dots,n\}$.

Now suppose $(v,w)\in V\times W$. There exist
$a_1,\dots,a_m,b_1,\dots,b_n\in\F$ such that $v=a_1e_1+\cdots+a_me_m$ and
$w=b_1f_1+\cdots+b_nf_n$. Thus
\begin{align*}
\hat\Gamma(v\otimes w)
  &= \hat\Gamma\biggl(\sum_{k=1}^n\sum_{j=1}^m (a_jb_k)(e_j\otimes f_k)\biggr)\\
  &= \sum_{k=1}^n\sum_{j=1}^m a_jb_k\hat\Gamma(e_j\otimes f_k)\\
  &= \sum_{k=1}^n\sum_{j=1}^m a_jb_k\Gamma(e_j,f_k)\\
  &= \Gamma(v,w),
\end{align*}
as desired, where the second line holds because $\hat\Gamma$ is linear, the
third line holds by the definition of $\hat\Gamma$, and the fourth line holds
because $\Gamma$ is bilinear.

The uniqueness of the linear map $\hat\Gamma$ satisfying
$\hat\Gamma(v\otimes w)=\Gamma(v,w)$ follows from \ref{ax:9.74}(b), completing
the proof of (a).

To prove (b), define a function $T^\#\colon V\times W\to U$ by
$T^\#(v,w)=T(v\otimes w)$ for all $(v,w)\in V\times W$. The bilinearity of the
tensor product (see \ref{ax:9.73}) and the linearity of $T$ imply that $T^\#$
is bilinear.

Clearly the choice of $T^\#$ that satisfies the conditions is unique.
\end{proof}

To prove \ref{ax:9.79}(a), we could not just define
$\hat\Gamma(v\otimes w)=\Gamma(v,w)$ for all $v\in V$ and $w\in W$ (and then
extend $\hat\Gamma$ linearly to all of $V\otimes W$) because elements of
$V\otimes W$ do not have unique representations as finite sums of elements of
the form $v\otimes w$. Our proof used a basis of $V$ and a basis of $W$ to get
around this problem.

Although our construction of $\hat\Gamma$ in the proof of \ref{ax:9.79}(a)
depended on a basis of $V$ and a basis of $W$, the equation
$\hat\Gamma(v\otimes w)=\Gamma(v,w)$ that holds for all $v\in V$ and $w\in W$
shows that $\hat\Gamma$ does not depend on the choice of bases for $V$ and $W$.

\subsection{Tensor Product of Inner Product Spaces}

The result below features three inner products---one on $V\otimes W$, one on
$V$, and one on $W$, although we use the same symbol
$\langle\cdot,\cdot\rangle$ for all three inner products.

\begin{result}[Inner product on tensor product of two inner product
spaces]\label{ax:9.80}
Suppose $V$ and $W$ are inner product spaces. Then there is a unique inner
product on $V\otimes W$ such that
\[ \ip{v\otimes w}{u\otimes x} = \ip{v}{u}\ip{w}{x} \]
for all $v,u\in V$ and $w,x\in W$.
\end{result}

\begin{proof}
Suppose $e_1,\dots,e_m$ is an orthonormal basis of $V$ and $f_1,\dots,f_n$ is
an orthonormal basis of $W$. Define an inner product on $V\otimes W$
by\axeqstep{ax:9.81}
\begin{equation*}
\biggl\langle\sum_{k=1}^n\sum_{j=1}^m b_{j,k}\,e_j\otimes f_k,\,
  \sum_{k=1}^n\sum_{j=1}^m c_{j,k}\,e_j\otimes f_k\biggr\rangle
  = \sum_{k=1}^n\sum_{j=1}^m b_{j,k}\overline{c_{j,k}}. \axeqtag
\end{equation*}
The straightforward verification that \ref{ax:9.81} defines an inner product
on $V\otimes W$ is left to the reader [use \ref{ax:9.74}(b)].

Suppose that $v,u\in V$ and $w,x\in W$. Let $v_1,\dots,v_m\in\F$ be such that
$v=v_1e_1+\cdots+v_me_m$, with similar expressions for $u$, $w$, and $x$. Then
\begin{align*}
\ip{v\otimes w}{u\otimes x}
  &= \biggl\langle\sum_{j=1}^m v_je_j\otimes\sum_{k=1}^n w_kf_k,\,
     \sum_{j=1}^m u_je_j\otimes\sum_{k=1}^n x_kf_k\biggr\rangle\\
  &= \biggl\langle\sum_{k=1}^n\sum_{j=1}^m v_jw_k\,e_j\otimes f_k,\,
     \sum_{k=1}^n\sum_{j=1}^m u_jx_k\,e_j\otimes f_k\biggr\rangle\\
  &= \sum_{k=1}^n\sum_{j=1}^m v_j\overline{u_j}w_k\overline{x_k}\\
  &= \biggl(\sum_{j=1}^m v_j\overline{u_j}\biggr)
     \biggl(\sum_{k=1}^n w_k\overline{x_k}\biggr)\\
  &= \ip{v}{u}\ip{w}{x}.
\end{align*}

There is only one inner product on $V\otimes W$ such that
$\ip{v\otimes w}{u\otimes x}=\ip{v}{u}\ip{w}{x}$ for all $v,u\in V$ and
$w,x\in W$ because every element of $V\otimes W$ can be written as a linear
combination of elements of the form $v\otimes w$ [by \ref{ax:9.74}(b)].
\end{proof}

The definition below of a natural inner product on $V\otimes W$ is now
justified by \ref{ax:9.80}. We could not have simply defined
$\ip{v\otimes w}{u\otimes x}$ to be $\ip{v}{u}\ip{w}{x}$ (and then used
additivity in each slot separately to extend the definition to $V\otimes W$)
without some proof because elements of $V\otimes W$ do not have unique
representations as finite sums of elements of the form $v\otimes w$.

\begin{definition}[Inner product on tensor product of two inner product
spaces]\label{ax:9.82}
Suppose $V$ and $W$ are inner product spaces. The inner product on $V\otimes W$
is the unique function $\langle\cdot,\cdot\rangle$ from
$(V\otimes W)\times(V\otimes W)$ to $\F$ such that
\[ \ip{v\otimes w}{u\otimes x} = \ip{v}{u}\ip{w}{x} \]
for all $v,u\in V$ and $w,x\in W$.
\end{definition}

Take $u=v$ and $x=w$ in the equation above and then take square roots to show
that
\[ \norm{v\otimes w} = \norm{v}\,\norm{w} \]
for all $v\in V$ and all $w\in W$.

The construction of the inner product in the proof of \ref{ax:9.80} depended
on an orthonormal basis $e_1,\dots,e_m$ of $V$ and an orthonormal basis
$f_1,\dots,f_n$ of $W$. Formula \ref{ax:9.81} implies that
$\{e_j\otimes f_k\}_{j=1,\dots,m;\,k=1,\dots,n}$ is a doubly indexed
orthonormal list in $V\otimes W$ and hence is an orthonormal basis of
$V\otimes W$ [by \ref{ax:9.74}(b)]. The importance of the next result arises
because the orthonormal bases used there can be different from the orthonormal
bases used to define the inner product in \ref{ax:9.80}. Although the notation
for the bases is the same in the proof of \ref{ax:9.80} and in the result
below, think of them as two different sets of orthonormal bases.

\begin{result}[Orthonormal basis of $V\otimes W$]\label{ax:9.83}
Suppose $V$ and $W$ are inner product spaces, and $e_1,\dots,e_m$ is an
orthonormal basis of $V$ and $f_1,\dots,f_n$ is an orthonormal basis of $W$.
Then
\[ \{e_j\otimes f_k\}_{j=1,\dots,m;\,k=1,\dots,n} \]
is an orthonormal basis of $V\otimes W$.
\end{result}

\begin{proof}
We know that $\{e_j\otimes f_k\}_{j=1,\dots,m;\,k=1,\dots,n}$ is a basis of
$V\otimes W$ [by \ref{ax:9.74}(b)]. Thus we only need to verify
orthonormality. To do this, suppose $j,M\in\{1,\dots,m\}$ and
$k,N\in\{1,\dots,n\}$. Then
\[ \ip{e_j\otimes f_k}{e_M\otimes f_N} = \ip{e_j}{e_M}\ip{f_k}{f_N}
   = \begin{cases} 1 & \text{if } j=M \text{ and } k=N,\\
                   0 & \text{otherwise.} \end{cases} \]
Hence the doubly indexed list $\{e_j\otimes f_k\}_{j=1,\dots,m;\,k=1,\dots,n}$
is indeed an orthonormal basis of $V\otimes W$.
\end{proof}

See Exercise 11 for an example of how the inner product structure on
$V\otimes W$ interacts with operators on $V$ and $W$.

\subsection{Tensor Product of Multiple Vector Spaces}

We have been discussing properties of the tensor product of two
finite-dimensional vector spaces. Now we turn our attention to the tensor
product of multiple finite-dimensional vector spaces. This generalization
requires no new ideas, only some slightly more complicated notation. Readers
with a good understanding of the tensor product of two vector spaces should be
able to make the extension to the tensor product of more than two vector
spaces.

Thus in this subsection, no proofs will be provided. The definitions and the
statements of results that will be provided should be enough information to
enable readers to fill in the details, using what has already been learned
about the tensor product of two vector spaces.

We begin with the following notational assumption.

\begin{notation}[$V_1,\dots,V_m$]\label{ax:9.84}
For the rest of this subsection, $m$ denotes an integer greater than $1$ and
$V_1,\dots,V_m$ denote finite-dimensional vector spaces.
\end{notation}

The notion of an $m$-linear functional, which we are about to define,
generalizes the notion of a bilinear functional (see \ref{ax:9.68}). Recall
that the use of the word ``functional'' indicates that we are mapping into the
scalar field $\F$. Recall also that the terminology ``$m$-linear form'' is used
in the special case $V_1=\cdots=V_m$ (see \ref{ax:9.25}). The notation
$\mathcal{B}(V_1,\dots,V_m)$ generalizes our previous notation
$\mathcal{B}(V,W)$.

\begin{definition}[$m$-linear functional, the vector space
$\mathcal{B}(V_1,\dots,V_m)$]\label{ax:9.85}\leavevmode\emergencystretch=3em
\begin{itemize}
\item An \emph{$m$-linear functional} on $V_1\times\cdots\times V_m$ is a
  function $\beta\colon V_1\times\cdots\times V_m\to\F$ that is a linear
  functional in each slot when the other slots are held fixed.
\item The vector space of $m$-linear functionals on $V_1\times\cdots\times V_m$
  is denoted by $\mathcal{B}(V_1,\dots,V_m)$.
\end{itemize}
\end{definition}

\begin{example}[$m$-linear functional]\label{ax:9.86}\emergencystretch=3em
Suppose $\varphi_k\in V_k{}'$ for each $k\in\{1,\dots,m\}$. Define
$\beta\colon V_1\times\cdots\times V_m\to\F$ by
\[ \beta(v_1,\dots,v_m) = \varphi_1(v_1)\times\cdots\times\varphi_m(v_m). \]
Then $\beta$ is an $m$-linear functional on $V_1\times\cdots\times V_m$.
\end{example}

The next result can be proved by imitating the proof of \ref{ax:9.70}.

\begin{result}[Dimension of the vector space of $m$-linear
functionals]\label{ax:9.87}\leavevmode
\[ \dim\mathcal{B}(V_1,\dots,V_m)=(\dim V_1)\times\cdots\times(\dim V_m). \]
\end{result}

Now we can define the tensor product of multiple vector spaces and the tensor
product of elements of those vector spaces. The following definition is
completely analogous to our previous definition (\ref{ax:9.71}) in the case
$m=2$.

\begin{definition}[Tensor product, $V_1\otimes\cdots\otimes V_m$,
$v_1\otimes\cdots\otimes v_m$]\label{ax:9.88}\leavevmode
\begin{itemize}
\item The \emph{tensor product} $V_1\otimes\cdots\otimes V_m$ is defined to be
  $\mathcal{B}(V_1',\dots,V_m')$.
\item For $v_1\in V_1,\dots,v_m\in V_m$, the \emph{tensor product}
  $v_1\otimes\cdots\otimes v_m$ is the element of $V_1\otimes\cdots\otimes V_m$
  defined by
  \[ (v_1\otimes\cdots\otimes v_m)(\varphi_1,\dots,\varphi_m)
     = \varphi_1(v_1)\cdots\varphi_m(v_m) \]
  for all $(\varphi_1,\dots,\varphi_m)\in V_1'\times\cdots\times V_m'$.
\end{itemize}
\end{definition}

The next result can be proved by following the pattern of the proof of the
analogous result when $m=2$ (see \ref{ax:9.72}).

\begin{result}[Dimension of the tensor product]\label{ax:9.89}\leavevmode
\[ \dim(V_1\otimes\cdots\otimes V_m)=(\dim V_1)\cdots(\dim V_m). \]
\end{result}

Our next result generalizes \ref{ax:9.74}.

\begin{result}[Basis of $V_1\otimes\cdots\otimes V_m$]\label{ax:9.90}
Suppose $\dim V_k=n_k$ and $e^k_1,\dots,e^k_{n_k}$ is a basis of $V_k$ for
$k=1,\dots,m$. Then
\[ \{e^1_{j_1}\otimes\cdots\otimes e^m_{j_m}\}_{j_1=1,\dots,n_1;\,\cdots;\,
   j_m=1,\dots,n_m} \]
is a basis of $V_1\otimes\cdots\otimes V_m$.
\end{result}

Suppose $m=2$ and $e^1_1,\dots,e^1_{n_1}$ is a basis of $V_1$ and
$e^2_1,\dots,e^2_{n_2}$ is a basis of $V_2$. Then with respect to the basis
$\{e^1_{j_1}\otimes e^2_{j_2}\}_{j_1=1,\dots,n_1;\,j_2=1,\dots,n_2}$ in the
result above, the coefficients of an element of $V_1\otimes V_2$ can be
represented by an $n_1$-by-$n_2$ matrix that contains the coefficient of
$e^1_{j_1}\otimes e^2_{j_2}$ in row $j_1$, column $j_2$. Thus we need a matrix,
which is an array specified by two indices, to represent an element of
$V_1\otimes V_2$.

If $m>2$, then the result above shows that we need an array specified by $m$
indices to represent an arbitrary element of $V_1\otimes\cdots\otimes V_m$.
Thus tensor products may appear when we deal with objects specified by arrays
with multiple indices.

The next definition generalizes the notion of a bilinear map (see
\ref{ax:9.77}). As with bilinear maps, the target space can be an arbitrary
vector space.

\begin{definition}[$m$-linear map]\label{ax:9.91}
An \emph{$m$-linear map} from $V_1\times\cdots\times V_m$ to a vector space $U$
is a function $\Gamma\colon V_1\times\cdots\times V_m\to U$ that is a linear
map in each slot when the other slots are held fixed.
\end{definition}

The next result can be proved by following the pattern of the proof of
\ref{ax:9.79}.

\begin{result}[Converting $m$-linear maps to linear maps]\label{ax:9.92}
Suppose $U$ is a vector space.
\begin{parts}
\item Suppose that $\Gamma\colon V_1\times\cdots\times V_m\to U$ is an
  $m$-linear map. Then there exists a unique linear map
  $\hat\Gamma\colon V_1\otimes\cdots\otimes V_m\to U$ such that
  \[ \hat\Gamma(v_1\otimes\cdots\otimes v_m) = \Gamma(v_1,\dots,v_m) \]
  for all $(v_1,\dots,v_m)\in V_1\times\cdots\times V_m$.
\item Conversely, suppose $T\colon V_1\otimes\cdots\otimes V_m\to U$ is a
  linear map. Then there exists a unique $m$-linear map
  $T^\#\colon V_1\times\cdots\times V_m\to U$ such that
  \[ T^\#(v_1,\dots,v_m) = T(v_1\otimes\cdots\otimes v_m) \]
  for all $(v_1,\dots,v_m)\in V_1\times\cdots\times V_m$.
\end{parts}
\end{result}

See Exercises 12 and 13 for tensor products of multiple inner product spaces.

% ------------------------------------------------------------------ exercises
\begin{exc}

\begin{exercise}
Suppose $v\in V$ and $w\in W$. Prove that $v\otimes w=0$ if and only if $v=0$
or $w=0$.
\end{exercise}

\begin{exercise}
Give an example of six distinct vectors $v_1,v_2,v_3,w_1,w_2,w_3$ in $\R^3$
such that
\[ v_1\otimes w_1+v_2\otimes w_2+v_3\otimes w_3 = 0 \]
but none of $v_1\otimes w_1,v_2\otimes w_2,v_3\otimes w_3$ is a scalar multiple
of another element of this list.
\end{exercise}

\begin{exercise}
Suppose that $v_1,\dots,v_m$ is a linearly independent list in $V$. Suppose
also that $w_1,\dots,w_m$ is a list in $W$ such that
\[ v_1\otimes w_1+\cdots+v_m\otimes w_m = 0. \]
Prove that $w_1=\cdots=w_m=0$.
\end{exercise}

\begin{exercise}
Suppose $\dim V>1$ and $\dim W>1$. Prove that
\[ \{v\otimes w : (v,w)\in V\times W\} \]
is not a subspace of $V\otimes W$.
\exnote{This exercise implies that if $\dim V>1$ and $\dim W>1$, then
\[ \{v\otimes w : (v,w)\in V\times W\} \ne V\otimes W. \]}
\end{exercise}

\begin{exercise}
Suppose $m$ and $n$ are positive integers. For $v\in\F^m$ and $w\in\F^n$,
identify $v\otimes w$ with an $m$-by-$n$ matrix as in Example~\ref{ax:9.76}.
With that identification, show that the set
\[ \{v\otimes w : v\in\F^m \text{ and } w\in\F^n\} \]
is the set of $m$-by-$n$ matrices (with entries in $\F$) that have rank at most
one.
\end{exercise}

\begin{exercise}
Suppose $m$ and $n$ are positive integers. Give a description, analogous to
Exercise 5, of the set of $m$-by-$n$ matrices (with entries in $\F$) that have
rank at most two.
\end{exercise}

\begin{exercise}
Suppose $\dim V>2$ and $\dim W>2$. Prove that
\[ \{v_1\otimes w_1+v_2\otimes w_2 : v_1,v_2\in V \text{ and } w_1,w_2\in W\}
   \ne V\otimes W. \]
\end{exercise}

\begin{exercise}
Suppose $v_1,\dots,v_m\in V$ and $w_1,\dots,w_m\in W$ are such that
\[ v_1\otimes w_1+\cdots+v_m\otimes w_m = 0. \]
Suppose that $U$ is a vector space and $\Gamma\colon V\times W\to U$ is a
bilinear map. Show that
\[ \Gamma(v_1,w_1)+\cdots+\Gamma(v_m,w_m) = 0. \]
\end{exercise}

\begin{exercise}
Suppose $S\in\Lin(V)$ and $T\in\Lin(W)$. Prove that there exists a unique
operator on $V\otimes W$ that takes $v\otimes w$ to $Sv\otimes Tw$ for all
$v\in V$ and $w\in W$.
\exnote{In an abuse of notation, the operator on $V\otimes W$ given by this
exercise is often called $S\otimes T$.}
\end{exercise}

\begin{exercise}
Suppose $S\in\Lin(V)$ and $T\in\Lin(W)$. Prove that $S\otimes T$ is an
invertible operator on $V\otimes W$ if and only if both $S$ and $T$ are
invertible operators. Also, prove that if both $S$ and $T$ are invertible
operators, then $(S\otimes T)^{-1}=S^{-1}\otimes T^{-1}$, where we are using the
notation from the comment after Exercise 9.
\end{exercise}

\begin{exercise}
Suppose $V$ and $W$ are inner product spaces. Prove that if $S\in\Lin(V)$ and
$T\in\Lin(W)$, then $(S\otimes T)^*=S^*\otimes T^*$, where we are using the
notation from the comment after Exercise 9.
\end{exercise}

\begin{exercise}
Suppose that $V_1,\dots,V_m$ are finite-dimensional inner product spaces.
Prove that there is a unique inner product on $V_1\otimes\cdots\otimes V_m$
such that
\[ \ip{v_1\otimes\cdots\otimes v_m}{u_1\otimes\cdots\otimes u_m}
   = \ip{v_1}{u_1}\cdots\ip{v_m}{u_m} \]
for all $(v_1,\dots,v_m)$ and $(u_1,\dots,u_m)$ in $V_1\times\cdots\times V_m$.
\exnote{Note that the equation above implies that
\[ \norm{v_1\otimes\cdots\otimes v_m} = \norm{v_1}\times\cdots\times\norm{v_m} \]
for all $(v_1,\dots,v_m)\in V_1\times\cdots\times V_m$.}
\end{exercise}

\begin{exercise}
Suppose that $V_1,\dots,V_m$ are finite-dimensional inner product spaces and
$V_1\otimes\cdots\otimes V_m$ is made into an inner product space using the
inner product from Exercise 12. Suppose $e^k_1,\dots,e^k_{n_k}$ is an
orthonormal basis of $V_k$ for each $k=1,\dots,m$. Show that the list
\[ \{e^1_{j_1}\otimes\cdots\otimes e^m_{j_m}\}_{j_1=1,\dots,n_1;\,\cdots;\,
   j_m=1,\dots,n_m} \]
is an orthonormal basis of $V_1\otimes\cdots\otimes V_m$.
\end{exercise}
\end{exc}
